\documentclass[11pt]{article}

\usepackage[T1]{fontenc}
\usepackage{lmodern}
\usepackage{amsmath,amssymb,amsthm}
\usepackage[a4paper,margin=25mm]{geometry}
\usepackage[hidelinks]{hyperref}

\newtheorem{theorem}{Theorem}
\newtheorem{lemma}[theorem]{Lemma}
\newtheorem{proposition}[theorem]{Proposition}
\newtheorem{conjecture}[theorem]{Conjecture}
\theoremstyle{definition}
\newtheorem{definition}[theorem]{Definition}
\newtheorem{example}[theorem]{Example}
\theoremstyle{remark}
\newtheorem*{status}{Status}

\title{<Thread key>: <Mathematics note title>}
\author{}
\date{\today}

\begin{document}
\maketitle

\begin{status}
State whether this is a definition note, example, proof attempt, checked proof, or counterexample.
Name the Question or Claim key and every Source-register ID with a precise locator.
\end{status}

\section{Question and assumptions}

State the question, hypotheses, notation, and scope.

\section{Definitions and examples}

Write each definition in your own notation, then test it on a small example.

\section{Argument}

Write the proof or proof attempt. Mark every unresolved step explicitly.

\section{Checks and dependencies}

List the results used, where they are established, and the examples or computations that test the
argument. A computation is evidence for behavior, not a proof of a general statement.

\section{What I can reconstruct}

State what you can reproduce without looking and what remains open.

\end{document}
